% =====================================================================
\section{ESPUI: веб-панель из программы C++}\label{les:espui}
% =====================================================================
\lessonintro
  {собрать веб-интерфейс из готовых элементов и связать его с состоянием устройства}
  {уроки 5 (Wi-Fi и первые панели ESPUI), 8 (очереди), 18 (JSON), 38 (надёжность)}
  {панель с переключателем LED, ползунком яркости, вкладками и показаниями АЦП}
  {ESPUI и её зависимости; LED с резистором на GPIO4, потенциометр на GPIO1}

% Фирменные цвета ESPUI для рисунков главы
\definecolor{turq}{HTML}{1ABC9C}
\definecolor{emerald}{HTML}{2ECC71}
\definecolor{peter}{HTML}{3498DB}
\definecolor{wetasph}{HTML}{34495E}
\definecolor{sunfl}{HTML}{F1C40F}
\definecolor{carrot}{HTML}{E67E22}
\definecolor{alizar}{HTML}{E74C3C}
\definecolor{darkc}{HTML}{444444}

\subsection{От первых панелей к полному каталогу}

С ESPUI вы уже работали: в уроке 5 мы собрали панель из подписи, переключателя и ползунка, а в уроках 6, 15, 16, 17, 24, 25 и 26 добавляли к ней показания и органы управления. Здесь мы соберём всё в систему: каталог элементов, вкладки, одна модель состояния для кнопок и браузера и проверка надёжности. Если понадобится нестандартный дизайн или собственный протокол, пригодится ручная страница на HTML, CSS и JavaScript из урока 44.

Первую страницу браузер получает по HTTP, дальнейшие события и обновления идут через WebSocket.
Соединение двустороннее: браузер отправляет команду, плата может сообщить новое показание,
не дожидаясь перезагрузки страницы. Наличие страницы ещё не доказывает работу WebSocket:
если кнопки не действуют, проверяйте отдельно канал событий.

\subsection{Установка: проверяем комплект зависимостей}

Примеры главы скомпилированы для ESP32-S3 с Arduino-ESP32 3.2.0.
Проверяемые версии зависимостей приведены ниже; это комплект для воспроизводимого опыта,
а не требование всегда использовать только эти выпуски.

\begin{center}\small
\begin{tabularx}{\textwidth}{@{}l l X@{}}
\toprule
Библиотека & Версия & Назначение\\
\midrule
ESPUI & 2.2.4 & элементы интерфейса и связь с браузером\\
ESPAsyncWebServer (ESP32Async) & 3.12.1 & HTTP, WebSocket, обслуживание клиентов\\
AsyncTCP (ESP32Async) & 3.5.0 & асинхронный TCP для ESP32\\
ArduinoJson & 7.4.2 & представление сообщений\\
\bottomrule
\end{tabularx}
\end{center}

В менеджере Arduino IDE установите ESPUI, ArduinoJson и зависимости из семейства ESP32Async.
У сетевых библиотек встречаются старые копии с теми же заголовками: смотрите в журнале сборки,
какой каталог действительно выбран. Не держите несколько конкурирующих версий.
На ESP32-S3 нужен AsyncTCP, а не ESPAsyncTCP для ESP8266.
LittleFS уже входит в используемое ядро; отдельная старая LittleFS-библиотека здесь не требуется.
Исходники и точные Git-коммиты проверенного набора записаны в
\code{examples/web-panels/README.md}.

\subsection{Точка доступа и ID элемента}

Первую панель из подписи, переключателя и ползунка вы собрали в уроке 5 (шаги 1--3 подраздела
«ESPUI: панель управления без HTML»), поэтому здесь её не повторяем. Все примеры этой главы
тоже начинаются с точки доступа: ошибка пароля домашней сети не мешает изучению интерфейса.
Подключите телефон к сети, имя которой указано в скетче, с учебным паролем \code{lesson-39-demo}
и откройте адрес из Serial Monitor, например \code{http://192.168.4.1/}. Если телефон пишет
«без интернета», останьтесь в этой сети. Для собственной платы задайте индивидуальный пароль.

Напомним главное. У элемента есть \textbf{ID}, подпись, цвет, значение и при необходимости обработчик
вида \code{void callback(Control* sender, int event)}. ID не является номером GPIO: он обозначает
виджет, который может управлять любым объектом. Значение доступно как \code{sender->value},
то есть строка. Ресурсы страницы берутся из прошивки, загрузка папки \code{data} во Flash не нужна,
а вызывать \code{handleClient()}, как у обычного WebServer, не требуется: сетевой стек асинхронный.

Для действия «выключить» достаточно обработать одно событие кнопки, например отпускание.
Если выполнять переключение и на нажатии, и на отпускании, один щелчок вернёт исходное состояние.
Для удерживаемого управления нужна другая модель: команда старта, команда остановки и
тайм-аут на случай потери клиента. Отпускание в браузере не является аппаратной аварийной остановкой.

\subsection{Каталог виджетов}

Сначала посмотрим, что вообще умеет библиотека. Каждый виджет --- это либо \emph{вход}
(пользователь что-то сообщает плате), либо \emph{выход} (плата что-то показывает),
либо украшение. Заголовок и цвет у всех одинаковые, различаются значение и событие.

\begin{center}\small
\begin{tabularx}{\textwidth}{@{}l l X l@{}}
\toprule
Виджет & Тип события & Что приходит и как показывать & Для чего\\
\midrule
Label & --- & \code{updateLabel(id, text)} & число, состояние, ошибка\\
Button & \code{B\_DOWN}, \code{B\_UP} & нажали / отпустили; текст кнопки в \code{value} & действие один раз\\
Switcher & \code{S\_ACTIVE}, \code{S\_INACTIVE} & \code{updateSwitcher(id, on)} & вкл/выкл режима\\
Slider & \code{SL\_VALUE} & число в строке; границы задают \code{Min}, \code{Max} & яркость, скорость\\
Number & \code{N\_VALUE} & число; границы тоже \code{Min}, \code{Max} & угол, порог, период\\
Text & \code{T\_VALUE} & произвольная строка --- её надо проверять & имя, пароль, команда\\
Select & \code{S\_VALUE} & в \code{value} лежит значение выбранного \code{Option} & режим, нота, профиль\\
Pad / PadWithCenter & \code{P\_FOR\_DOWN}, \code{P\_FOR\_UP}, \ldots{} & пять кнопок со стрелками (с центром --- шесть) & движение, меню\\
Gauge & --- & \code{updateControlValue(id, строка)}, \code{Min}, \code{Max} & круговая шкала\\
Graph & --- & \code{addGraphPoint(id, число)}, \code{clearGraph(id)} & показания за время\\
Separator & --- & только подпись & разделить группы\\
Accelerometer & --- & датчик наклона самого телефона & управление наклоном\\
\bottomrule
\end{tabularx}
\end{center}

Цвет задаётся одним из восьми имён: \code{Turquoise}, \code{Emerald}, \code{Peterriver},
\code{Wetasphalt}, \code{Sunflower}, \code{Carrot}, \code{Alizarin}, \code{Dark}. Они не несут
смысла сами по себе, поэтому договоритесь заранее: например, красный (\code{Alizarin}) ---
только для опасных действий и ошибок, зелёный (\code{Emerald}) --- для включения,
синий (\code{Peterriver}) --- для чисел и измерений. Тогда по цвету видно, что делает элемент,
как на цветовой схеме книги.

\begin{figure}[H]
\centering
\begin{tikzpicture}[x=1cm,y=1cm,font=\sffamily\scriptsize]
  \tikzset{sw/.style={minimum width=18.5mm,minimum height=9mm,rounded corners=1.5pt,text=white,font=\sffamily\scriptsize\bfseries,align=center}}
  \node[sw,fill=turq] at (0,0) {Turquoise};
  \node[sw,fill=emerald] at (2,0) {Emerald};
  \node[sw,fill=peter] at (4,0) {Peterriver};
  \node[sw,fill=wetasph] at (6,0) {Wetasphalt};
  \node[sw,fill=sunfl,text=black] at (8,0) {Sunflower};
  \node[sw,fill=carrot] at (10,0) {Carrot};
  \node[sw,fill=alizar] at (12,0) {Alizarin};
  \node[sw,fill=darkc] at (14,0) {Dark};
  \node[anchor=north,text=gray] at (1,-0.6) {подпись};
  \node[anchor=north,text=okgreen!70!black] at (3,-0.6) {включение};
  \node[anchor=north,text=gray] at (5,-0.6) {числа};
  \node[anchor=north,text=gray] at (7,-0.6) {шкала};
  \node[anchor=north,text=gray] at (9,-0.6) {внимание};
  \node[anchor=north,text=gray] at (11,-0.6) {сервис};
  \node[anchor=north,text=red!70!black] at (13,-0.6) {стоп, ошибка};
\end{tikzpicture}
\caption{Восемь цветов ESPUI. Закрепите за ними роли, например: красный --- «стоп» и ошибки, зелёный --- включение, синий --- числа}
\end{figure}

\subsection{Как выглядит панель}

Программа создаёт элементы по порядку, и страница выстраивает их в колонку карточек.
Вкладки (\code{Tab}) становятся кнопками вверху, а у каждого элемента \code{parentControl}
говорит, на какой вкладке он живёт. Ниже --- схематичный вид панели из примера \code{S3-Widgets}:
слева вкладка «Input» с полями ввода, справа вкладка «Output» с показаниями.

\begin{figure}[H]
\centering
\begin{tikzpicture}[x=1cm,y=1cm,font=\sffamily\scriptsize]
  \tikzset{card/.style={rounded corners=2pt,draw=gray!60,fill=white,minimum width=44mm,align=left,inner sep=3pt},
           lab/.style={font=\sffamily\tiny,text=gray},
           tab/.style={minimum width=22mm,minimum height=6mm,font=\sffamily\scriptsize\bfseries,text=white,inner sep=1pt}}
  % Левый телефон: Input
  \begin{scope}
    \fill[black!80,rounded corners=4mm] (-2.9,-10.6) rectangle (2.9,1.2);
    \fill[gray!10] (-2.6,-10.2) rectangle (2.6,0.7);
    \node[tab,fill=peter] at (-1.3,0.3) {Input};
    \node[tab,fill=gray!55] at (1.3,0.3) {Output};
    \node[lab,anchor=west] at (-2.4,-0.6) {Ввод};
    \draw[gray!50] (-2.4,-0.85) -- (2.4,-0.85);
    \node[card,minimum height=11mm] at (0,-1.9) {};
    \node[lab,anchor=west] at (-2.3,-1.5) {Number 0..50};
    \node[draw=peter,fill=white,minimum width=40mm,minimum height=5mm,anchor=north west] at (-2.2,-1.7) {10};
    \node[card,minimum height=11mm] at (0,-3.6) {};
    \node[lab,anchor=west] at (-2.3,-3.2) {Text};
    \node[draw=emerald,fill=white,minimum width=40mm,minimum height=5mm,anchor=north west] at (-2.2,-3.4) {hello};
    \node[card,minimum height=11mm] at (0,-5.3) {};
    \node[lab,anchor=west] at (-2.3,-4.9) {Select};
    \node[draw=sunfl,fill=white,minimum width=40mm,minimum height=5mm,anchor=north west] at (-2.2,-5.1) {do (C4) \ding{116}};
    \node[card,minimum height=38mm] at (0,-8.2) {};
    \node[lab,anchor=west] at (-2.3,-6.9) {Pad};
    \foreach \dx/\dy/\s in {0/0.9/↑,0/-0.9/↓,-0.9/0/←,0.9/0/→}
      \node[circle,fill=carrot,text=white,minimum size=7mm,font=\sffamily\small] at (\dx,-8.4+\dy) {\s};
    \node[circle,fill=carrot!60,text=white,minimum size=6mm,font=\sffamily\tiny] at (0,-8.4) {OK};
  \end{scope}
  % Правый телефон: Output
  \begin{scope}[xshift=7cm]
    \fill[black!80,rounded corners=4mm] (-2.9,-10.6) rectangle (2.9,1.2);
    \fill[gray!10] (-2.6,-10.2) rectangle (2.6,0.7);
    \node[tab,fill=gray!55] at (-1.3,0.3) {Input};
    \node[tab,fill=peter] at (1.3,0.3) {Output};
    \node[card,minimum height=33mm] at (0,-2.5) {};
    \node[lab,anchor=west] at (-2.3,-1.1) {Gauge, \%};
    \draw[gray!30,line width=5pt] (-1.4,-3.8) arc (180:0:1.4);
    \draw[wetasph,line width=5pt] (-1.4,-3.8) arc (180:105:1.4);
    \node[font=\sffamily\small\bfseries] at (0,-3.55) {40};
    \node[card,minimum height=33mm] at (0,-6.1) {};
    \node[lab,anchor=west] at (-2.3,-4.7) {Graph};
    \draw[gray!40] (-2.1,-7.5) -- (2.1,-7.5);
    \draw[turq,thick] (-2.1,-7.1) -- (-1.5,-6.6) -- (-0.9,-6.9) -- (-0.3,-5.9) -- (0.3,-6.4) -- (0.9,-5.4) -- (1.5,-6.0) -- (2.1,-5.5);
    \node[card,minimum height=11mm] at (0,-8.8) {};
    \node[lab,anchor=west] at (-2.3,-8.4) {Label};
    \node[anchor=west,font=\sffamily\small,text=alizar] at (-2.2,-9.0) {t = 53 s};
  \end{scope}
\end{tikzpicture}
\caption{Схема панели \code{S3-Widgets}: у каждого элемента подпись сверху и значение под ней; вкладки переключаются нажатием}
\end{figure}

Размеры карточек и порядок задаёт библиотека, поэтому в интерфейсе не нужно думать о
координатах. Если на странице не хватает места, широкую карточку можно сделать функцией
\code{setPanelWide(id, true)}, а набор кнопок разложить в строку через \code{setVertical(id, false)}.

\subsection{Gauge и Graph: вольтметр на телефоне}

В уроке 6 мы читали потенциометр и выводили число в Serial Monitor. Покажем то же самое
стрелкой и графиком. Подключение прежнее: крайние выводы потенциометра к 3V3 и GND,
движок к GPIO1.

\codecap{ESPUI: вольтметр со шкалой и графиком}
\codeinput{l39_espui_voltmeter}

Что здесь нового:
\begin{itemize}[leftmargin=*]
  \item у \code{Gauge} границы шкалы создаются дочерними элементами \code{Min} и \code{Max};
        значение обновляется через \code{updateControlValue} --- это строка, поэтому число
        оборачиваем в \code{String()};
  \item точки \code{Graph} добавляет \code{addGraphPoint}; график хранит ограниченное число
        последних отсчётов. Это окно просмотра, а не журнал: историю сохраняйте отдельно;
  \item чтение АЦП усредняется по 16 замерам, как в уроке 6, и выполняется два раза в секунду.
        Страница не должна обновляться чаще, чем глазу нужно.
\end{itemize}

\begin{tip}
Если стрелка Gauge «упирается» в край, проверьте \code{Min} и \code{Max}: библиотека
не отсекает значения сама. Напряжение выше 3,3~В на вывод подавать нельзя ни при какой шкале.
\end{tip}

\subsection{Ввод: числа, текст, списки и крестовина}

Скетч \code{S3-Widgets} не управляет внешними деталями, он нужен, чтобы увидеть все входные
виджеты и их события. Нажимайте и смотрите Serial Monitor: для каждого события печатается строка.

Соберём его по шагам, добавляя по одному виджету.

\subsubsection{Шаг 1. Числовое и текстовое поле}

Начнём с двух самых простых полей ввода на вкладке «Input». Подключитесь к сети \code{S3-Widgets},
введите число и текст и посмотрите, что печатает монитор.

\codecap{Шаг 1. Number и Text}
\codeinput{l39_espui_widgets_s1}

\begin{itemize}[leftmargin=*]
  \item \code{ControlType::Tab} создаёт вкладку, а её ID передаётся как родитель каждому элементу, который
        должен на ней жить;
  \item у \code{Number} границы задают дочерние элементы \code{Min} и \code{Max}, а событие \code{N\_VALUE}
        приходит, когда пользователь закончил ввод;
  \item \textbf{Text хранит строку, а не число.} Поле можно очистить или вставить в него что угодно.
        Всё, что приходит из \code{Text}, проверяйте так же строго, как команды из Serial в уроке 2.
\end{itemize}

\subsubsection{Шаг 2. Список выбора}

\code{Select} нужен, когда допустимых значений мало: нота, режим, профиль. Пользователь не может ввести
ничего «лишнего», поэтому проверять такой ввод проще.

\codecap{Шаг 2. Добавляем Select}
\codeinput{l39_espui_widgets_s2}

\begin{itemize}[leftmargin=*]
  \item сначала создаём сам список, затем несколько \code{Option} с \code{parentControl}, равным его ID;
  \item значение, которое придёт в \code{sender->value}, --- второй аргумент \code{Option}
        (здесь \code{"do"}, \code{"re"}, \code{"mi"}), а не подпись на экране;
  \item событие списка называется \code{S\_VALUE}, не путайте с \code{S\_ACTIVE} переключателя.
\end{itemize}

\subsubsection{Шаг 3. Крестовина Pad}

Крестовина из пяти кнопок (с центром --- шести) нужна для движения и меню. Её событий много, поэтому
обработчик устроен как \code{switch}.

\codecap{Шаг 3. Добавляем Pad}
\codeinput{l39_espui_widgets_s3}

\begin{itemize}[leftmargin=*]
  \item \code{PadWithCenter} добавляет шестую кнопку в середине;
  \item \textbf{Pad} шлёт пары «нажали--отпустили». Если по нажатию мотор едет, то по отпусканию он
        обязан остановиться; при потере связи отпускание может не прийти, поэтому нужен таймер;
  \item \textbf{Number и Slider отдают число строкой}, поэтому \code{toInt()} вернёт 0 на мусоре.
        Для настоящих «нулей» проверяйте длину и символы, как это делает панель дальше в уроке.
\end{itemize}

\subsubsection{Шаг 4. Вкладка Output: Gauge, Graph и Label}

Теперь вторая вкладка с \emph{выходами}. Датчика нет: значение «пилой» от 0 до 100 создаёт сама программа,
чтобы было видно, как работают шкала, график и подпись. Это итоговый скетч \code{l39\_espui\_widgets}.

\codecap{Шаг 4. Вкладка Output с показаниями}
\codeinput[\footnotesize]{l39_espui_widgets}

\begin{itemize}[leftmargin=*]
  \item второй \code{Tab} создаёт вторую кнопку-вкладку; элементы попадают на неё через свой родитель;
  \item \code{Gauge}, \code{Graph} и \code{Label} --- только выходы: событий у них нет, плата обновляет их
        сама через \code{updateControlValue}, \code{addGraphPoint} и \code{updateLabel};
  \item обновление раз в секунду (по \code{millis()}, как в уроке 5) не мешает остальной работе платы.
\end{itemize}

\begin{fact}
Раскладка панели не хранится в Flash отдельно: библиотека строит её заново из ваших вызовов
\code{addControl} при каждом запуске. Поэтому порядок создания в \code{setup()} --- это порядок
элементов на странице.
\end{fact}

\subsection{Одна модель состояния для кнопок и браузера}

Не поручайте сетевому callback долгие измерения, запись Flash или управление мотором.
У асинхронного сервера обработчик может выполняться вне \code{loop()}.
В нашем примере callback проверяет вход и ставит короткую команду в очередь;
\code{loop()} единолично меняет состояние LED и GPIO. Это защищает модель приложения
от одновременного изменения разными источниками команд.

\begin{figure}[H]\centering
\begin{tikzpicture}[node distance=5mm,font=\sffamily\scriptsize]
  \node[blk,fill=blkmem] (browser) {Браузер\\виджеты};
  \node[blk,fill=blkrf,right=of browser] (cb) {WebSocket\\callback};
  \node[blk,fill=blkrtc,right=of cb] (queue) {Очередь\\команд};
  \node[blk,fill=blkcpu,right=of queue] (model) {loop()\\on, level};
  \node[blk,fill=blkper,right=of model] (led) {PWM\\LED};
  \draw[arr] (browser)--(cb);
  \draw[arr] (cb)--(queue);
  \draw[arr] (queue)--(model);
  \draw[arr] (model)--(led);
  \draw[arr,accent] (model.south)--++(0,-0.8)-| node[pos=.5,below,font=\scriptsize]{публикация фактически применённого состояния} (browser.south);
\end{tikzpicture}
\caption{Команда интерфейса и состояние устройства проходят через общую модель}
\end{figure}

Ползунок ограничен 0--100 и в браузере, и на плате. Строка сначала проверяется на цифры и длину:
одного \code{toInt()} недостаточно, иначе ошибочный текст можно принять за ноль.
В очередь копируется число, а не указатель \code{Control*}, срок жизни которого не стоит
переносить за пределы обработчика.

\subsection{Полная панель: LED, вкладки и телеметрия}

GPIO4 подключите к LED через резистор \SI{220}{\ohm}, катод --- к GND.
Края потенциометра соедините с 3V3 и GND, движок --- с GPIO1. Другие опыты с этими GPIO отключите.
После прошивки подключитесь к \code{S3-Panel}, пароль тот же учебный.
Панель имеет вкладки управления и измерений; никаких реальных градусов из потенциометра
она не выдумывает: отображаются милливольты АЦП. Соберём её по шагам, как схему на рисунке выше:
сначала простейший переключатель, потом очередь, ползунок, кнопку остановки и в конце вторую вкладку.

\subsubsection{Шаг 1. Вкладка и переключатель LED}

Начнём с одного элемента, как в уроке 5, но создадим его общей функцией \code{addControl}: так
можно указать вкладку-родителя. Callback только записывает желание в \code{std::atomic}, а LED
включает \code{loop()}.

\codecap{Шаг 1. Одна вкладка и переключатель}
\codeinput{l39_espui_panel_s1}

\begin{itemize}[leftmargin=*]
  \item \code{addControl(тип, подпись, значение, цвет, родитель, обработчик)} --- универсальный способ
        создать любой элемент; родитель \code{mainTab} помещает его на вкладку «Control»;
  \item \code{wantOn} --- общая переменная двух задач, поэтому она атомарная (см. урок 5);
  \item подключитесь к сети \code{S3-Panel}: переключатель должен зажигать и гасить LED на GPIO4.
\end{itemize}

\subsubsection{Шаг 2. Очередь команд}

Одной переменной хватает, пока команда одна. Но если за короткое время придут «включить», «изменить
яркость» и «всё выключить», переменная запомнит только последнее значение, а порядок потеряется. Очередь
хранит команды по порядку, и \code{loop()} разбирает их одну за другой.

\codecap{Шаг 2. Callback кладёт команду в очередь}
\codeinput{l39_espui_panel_s2}

\begin{itemize}[leftmargin=*]
  \item \code{Command} --- маленькая структура «что сделать и с каким числом»; в очередь копируется число,
        а не указатель \code{Control*};
  \item \code{xQueueSend(..., 0)} не ждёт: если очередь полна, команда не попадает в неё, а счётчик
        \code{dropped} растёт (очереди FreeRTOS --- урок 8);
  \item переменная \code{on} теперь принадлежит только \code{loop()}: больше никто её не трогает.
\end{itemize}

\subsubsection{Шаг 3. Ползунок яркости с проверкой}

Ползунок отдаёт значение строкой. Пока оно не проверено, в очередь его нельзя: одного \code{toInt()}
недостаточно, потому что мусорный текст превратится в ноль.

\codecap{Шаг 3. Ползунок и проверка строки}
\codeinput{l39_espui_panel_s3}

\begin{itemize}[leftmargin=*]
  \item \code{parsePercent()} проверяет длину и символы и возвращает $-1$ для всего, что не число
        от 0 до 100;
  \item \code{Kind::Level} --- второй вид команды; \code{loop()} учитывает его в одном \code{if};
  \item границы \code{Min} и \code{Max} задают шкалу в браузере, а проверка на плате защищает от поддельных
        сообщений: доверять только браузеру нельзя.
\end{itemize}

\subsubsection{Шаг 4. Кнопка All off и применённое состояние}

Теперь браузер должен видеть то, что \emph{реально} сделала плата. Раз в 500~мс \code{loop()} публикует
в панель текущие \code{on} и \code{level}, поэтому переключатель и ползунок сами «возвращаются» на место
после кнопки аварийного выключения.

\codecap{Шаг 4. Кнопка и публикация состояния}
\codeinput[\footnotesize]{l39_espui_panel_s4}

\begin{itemize}[leftmargin=*]
  \item кнопка обрабатывает только \code{B\_UP}, иначе один щелчок сработает дважды;
  \item \code{updateSwitcher}, \code{updateSlider} и \code{updateLabel} отправляют в браузер
        \emph{применённое} состояние, а не то, что пользователь просил;
  \item публикация по таймеру (500~мс), а не в каждом проходе \code{loop()}: сеть не нужно заваливать.
\end{itemize}

\subsubsection{Шаг 5. Вторая вкладка с измерениями}

В последнем шаге появляется вкладка «Measurements»: милливольты АЦП на GPIO1 и число потерянных команд.
Это итоговый скетч \code{l39\_espui\_panel}.

\codecap{Шаг 5. Две вкладки: управление и измерения}
\codeinput[\footnotesize]{l39_espui_panel}

\begin{itemize}[leftmargin=*]
  \item вторая вкладка \code{dataTab} создаётся так же, как первая; метки \code{voltageId} и \code{queueId}
        получают её своим родителем;
  \item \code{analogReadMilliVolts(1)} читает АЦП и сразу переводит в милливольты (урок 6);
  \item проверка при запуске теперь останавливает программу, если очередь или ШИМ не создались:
        лучше явная ошибка в мониторе, чем панель, которая «молча» ничего не делает.
\end{itemize}

\subsection{Что проверять после запуска}

Включите LED и измените яркость, затем нажмите All off: переключатель и физический выход
должны показывать один результат. Откройте страницу во втором браузере.
Значения публикуются из общей модели каждые 500 мс, поэтому это не два независимых устройства.
При одновременных командах побеждает последняя обработанная команда; права владельца
или блокировку второго оператора при необходимости задаёт приложение.

\code{sliderContinuous=false} уменьшает поток событий: значение отправляется после окончания
перемещения. При непрерывном режиме задайте ограничение частоты и правило обработки переполнения.
Очередь примера конечна: потерянные команды считаются и видны на вкладке измерений.
Периодическая публикация возвращает браузеру применённое состояние, но не превращает
потерянную команду в выполненную. Для ответственных операций добавляют номер запроса и подтверждение.

Обновление 2 раза в секунду подходит для учебного АЦП. Не отправляйте новый JSON в каждом проходе
\code{loop()}: это тратит память и канал. Для датчика температуры обычно достаточно более редких
обновлений. Сохранение настроек в NVS делайте отдельной кнопкой «Сохранить» или после паузы,
а не на каждое движение ползунка. Подпись состояния и признак отсутствующего датчика полезнее
ложного нуля.

\subsection{Оформление, файлы и доступ}

Вкладки и группы создаются через \code{addControl} с ID родителя.
Для небольших изменений оформления есть \code{setPanelStyle} и \code{setElementStyle};
внешний вид настраивают после того, как работают значения, ошибки и синхронизация.
График --- виджет интерфейса, не долговременный журнал: историю измерений храните отдельно,
например на SD, в файле или внешней системе.

Для ресурсов из файловой системы есть \code{beginLITTLEFS()}.
Подготовка раздела --- отдельная операция: \code{prepareFileSystem()} по умолчанию допускает
форматирование. Не вставляйте её бездумно в каждый старт и не проверяйте на разделе с важными файлами.
В учебных скетчах главы используется обычный \code{begin()}, файловая система не меняется.

Пароль точки доступа ограничивает подключение к сети, но не является полной авторизацией команд.
Логин страницы также не доказывает защиту всех каналов: в проверенном ESPUI 2.2.4
флаг \code{WS\_AUTHENTICATION} по умолчанию выключен, а установка аутентификации WebSocket
зависит от этого флага. Перед применением в изделии отдельно проверьте HTTP, WebSocket,
шифрование и права команд. Примеры главы предназначены для локального учебного стенда.

\subsection{Те же панели в других уроках}

ESPUI особенно хорош как «пульт» для опытов, которые вы уже собрали. В нескольких уроках
книги есть врезка «То же на ESPUI»: к знакомой схеме добавляется телефон, а программа
датчика или мотора остаётся прежней. Все скетчи следуют правилу этого урока:
callback только записывает просьбу, а железо трогает \code{loop()}.

\begin{center}\small
\begin{tabularx}{\textwidth}{@{}l l X@{}}
\toprule
Урок & Скетч & Что показывает или чем управляет\\
\midrule
4 (ШИМ) & \code{l04\_espui\_gamma} & Slider яркости, Switcher гаммы, Label с коэффициентом заполнения\\
6 (АЦП) & \code{l39\_espui\_voltmeter} & Gauge и Graph напряжения потенциометра\\
15 (датчики) & \code{l15\_espui\_radar} & Gauge и Graph расстояния HC-SR04, сообщение «нет эха»\\
16 (мощная нагрузка) & \code{l16\_espui\_motor} & Slider скорости и Pad «вперёд/назад» с ограничением времени пуска\\
17 (адресные светодиоды) & \code{l17\_espui\_hue} & два Slider: оттенок и яркость встроенного RGB\\
24 (звук) & \code{l24\_espui\_notes} & Select ноты и Button «Играть»\\
25 (шаговый двигатель) & \code{l25\_espui\_stepper} & Number шага, кнопки «вправо/влево/в начало», Label угла\\
26 (акселерометр) & \code{l26\_espui\_tilt} & Label с «сырым» и фильтрованным углом, Graph наклона\\
\bottomrule
\end{tabularx}
\end{center}

\subsection{Какие ещё средства стоит знать}

\begin{center}\small
\begin{tabularx}{\textwidth}{@{}l X X@{}}
\toprule
Средство & Подходящая учебная задача & Что проверить перед выбором\\
\midrule
ESPUI & управление и диагностика готовыми виджетами & зависимости, поток событий, доступ\\
ESP-DASH & панель показаний, карточки и графики & различия открытой и Pro-версий\\
GyverLibs Settings & страница настроек, привязка данных и сохранение & комплект зависимостей, выбранный сервер\\
ESPAsyncWebServer & собственный HTML/CSS/JS и протокол & самостоятельно реализуем состояние и ввод\\
WiFiManager & первичная настройка Wi-Fi & это портал настройки, а не замена панели управления\\
\bottomrule
\end{tabularx}
\end{center}

ESP-DASH и Settings дают другой набор готовых компонентов; сравнивайте их на одной и той же
маленькой задаче, а не по числу картинок в README. В открытой и коммерческой редакциях могут
отличаться возможности и условия использования. Зависимости и API фиксируйте вместе со скетчем.
Эти альтернативы здесь рассмотрены для выбора; их отдельные программы не объявляются проверенными.

Если страницу пишете сами, выбирайте канал по задаче: обычный запрос удобен для разовой операции,
SSE --- для событий от сервера к браузеру, WebSocket --- для двустороннего обмена.
Любой канал требует понятного поведения при потере связи. WiFiManager разберём следующим уроком.

\begin{trouble}
\begin{itemize}[leftmargin=*]
  \item Не компилируется сеть --- проверьте версии, автора зависимостей и выбранные каталоги.
  \item Страница есть, кнопки молчат --- проверьте WebSocket, события и сохранённые ID.
  \item Одно нажатие действует дважды --- обработаны и B\_DOWN, и B\_UP.
  \item Значение «отскакивает» --- проверьте очередь, валидацию и источник истинного состояния.
  \item Телефон уходит из AP --- подтвердите работу без интернета и откройте адрес вручную.
\end{itemize}
\end{trouble}

\begin{summary}
ESPUI ускоряет создание типовой панели. ID обозначает виджет, callback сообщает событие,
а модель приложения хранит применённое состояние. Обновления, очередь, сохранение и
обработка разрыва связи остаются частью программы устройства.
\end{summary}

\begin{selfcheck}
\begin{enumerate}
  \item Чем ID элемента ESPUI отличается от номера GPIO?
  \item Почему долгое действие лучше вынести из callback в loop() или отдельную задачу?
  \item Почему успешная загрузка страницы и пароль HTTP ещё не доказывают защиту команд?
\end{enumerate}
\end{selfcheck}

\needspace{20\baselineskip}
\begin{tasks}
\begin{enumerate}
  \item \lv{1} Запустите первую панель из урока 5, измените заголовок и найдите ID подписи.
  \item \lv{1} Подключите панель и подтвердите, что All off выключает LED и переключатель.
  \item \lv{2} Добавьте кнопку установки яркости 25 процентов через ту же очередь команд.
  \item \lv{2} Подключите датчик температуры из урока 15; показывайте ошибку чтения вместо нуля.
  \item \lv{3} Сохраните настройки по отдельной кнопке, загрузите после перезапуска и проверьте два браузера.
  \item \lv{4} Проверьте потерю связи, переполнение очереди и повтор команд; добавьте номер команды, подтверждение и безопасное состояние для выбранного устройства.
\end{enumerate}
\end{tasks}

\paragraph{Документация для продолжения}
\href{https://github.com/s00500/ESPUI}{ESPUI: API и примеры},
\href{https://github.com/ESP32Async/ESPAsyncWebServer}{ESP32Async},
\href{https://github.com/ayushsharma82/ESP-DASH}{ESP-DASH},
\href{https://github.com/GyverLibs/Settings}{Settings}.
